\clearpage
\thispagestyle{plain}
\pagestyle{fancy}

\vspace*{1cm}
{\Huge\bfseries\raggedleft Preface: Personal motivation \par}
\vspace{1.5cm}
\addcontentsline{toc}{chapter}{Preface: Personal motivation}
\markboth{Preface: Personal motivation}{}

\section*{A personal journey}
From an early age, I was fascinated by characters who used science to understand and transform their world. Fictional figures like Tails, Bulma, and Iron Man, whose mastery of applied physics exemplifies the transformative power of knowledge, profoundly shaped my perspective. This early exposure sparked a lasting curiosity for inquiry and the systematic pursuit of understanding, which ultimately evolved into a genuine motivation for scientific research and my study of physics and chemistry.

\section*{Scientific motivation}
This dissertation is born from the conviction that physics (mathematics being the language to read the world\footnote{Galilei, G. \textit{Il Saggiatore}, Giacomo Mascardi, Roma, 1623.}, and physics its translation into tangible reality\footnote{Cecilia, C. E. ``Feynman Tells Difference Between Physics and Math'', \textit{The Cornell Daily Sun}, 1964, LXXXI(40).}) provides an essential vantage point to study the complexity of antimicrobial peptides. Rather than seeking definitive solutions, this work aims to clarify the underlying mechanisms of these interactions, offering a modest contribution to the global fight against antimicrobial resistance. It explores how the synergy between programming, molecular dynamics, and the emerging frontier of quantum computing can serve as a digital microscope, simulating the complex interactions of life at the atomic scale. Ultimately, its contribution lies not only in charting new methodologies, but also in sharing our unfruitful avenues, sparing future researchers from repeating the same mistakes.
